% Why a refactor of def_3_1 (CDMG) -- removing or restratifying the
% disjoint_EL field -- is the only LN-faithful way to close the
% 4-deviation marginalization cascade, what the failed attempt at
% def_3_14 taught us, and why the obstruction was found so late.
% Sibling article to 01_disjoint_EL_cascade.tex.
% Compiles via:
%   latexmk -pdflua 02_CDMG_disjoint_EL_refactor_needed.tex
% !TEX program = lualatex

\documentclass[11pt,a4paper]{article}

\usepackage[utf8]{inputenc}
\usepackage[T1]{fontenc}
\usepackage{amsmath,amssymb,amsthm}
\usepackage{stmaryrd}
\usepackage{xcolor}
\usepackage{tikz}
\usetikzlibrary{arrows.meta,decorations.markings}
\usepackage{hyperref}
\hypersetup{colorlinks=true,linkcolor=blue,urlcolor=blue}
\usepackage{listings}
\usepackage{booktabs}
\usepackage{pifont}  % for \ding (checkmarks + crosses in the option table)

% LN graph-notation macros (subset of leanification/preamble.tex).
\providecommand{\arrhead}{{Latex}}
\providecommand{\arrtail}{{}}
\providecommand{\tuh}[1][]{\mathrel{\tikz [baseline=-0.25ex,\arrtail-\arrhead, #1] \draw [#1] (0pt,0.5ex) -- (1.3em,0.5ex);}}
\providecommand{\hut}[1][]{\mathrel{\tikz [baseline=-0.25ex,\arrhead-\arrtail, #1] \draw [#1] (0pt,0.5ex) -- (1.3em,0.5ex);}}
\providecommand{\huh}[1][]{\mathrel{\tikz [baseline=-0.25ex,\arrhead-\arrhead, #1] \draw [#1] (0pt,0.5ex) -- (1.3em,0.5ex);}}

% Set / operator macros.
\providecommand{\sm}{\setminus}
\providecommand{\ins}{\subseteq}

% Independence / separation notation.
\providecommand{\Perp}{\operatorname{\perp}}
\providecommand{\nPerp}{\operatorname{\not\perp}}
\providecommand{\isPerp}{\Perp^{i\sigma}}
\providecommand{\nisPerp}{\nPerp^{i\sigma}}

% Graph predicates.
\providecommand{\Anc}{\mathrm{Anc}}

% Light Lean syntax highlighting (same as the sibling article).
\lstdefinelanguage{Lean}{
  morekeywords={def,theorem,lemma,structure,where,abbrev,by,intro,exact,
                refine,rintro,obtain,rcases,subst,rw,simp,Set,Prop,Type,
                Disjoint,Walk,CDMG,Fin,True,False,not,sorry,if,then,else,
                fun,let,in},
  sensitive=true,
  morecomment=[l]{--},
  morecomment=[s]{/-}{-/},
  morestring=[b]",
}
\lstset{
  language=Lean,
  basicstyle=\ttfamily\footnotesize,
  keywordstyle=\bfseries,
  commentstyle=\itshape\color{gray!70!black},
  breaklines=true,
  showstringspaces=false,
  columns=fullflexible,
  upquote=true,
  literate=
    {α}{{$\alpha$}}1 {β}{{$\beta$}}1 {π}{{$\pi$}}1 {σ}{{$\sigma$}}1
    {ℕ}{{$\mathbb{N}$}}1 {⊆}{{$\subseteq$}}1 {∈}{{$\in$}}1
    {∉}{{$\notin$}}1 {∅}{{$\emptyset$}}1 {∪}{{$\cup$}}1
    {∩}{{$\cap$}}1 {∧}{{$\wedge$}}1 {∨}{{$\vee$}}1 {¬}{{$\neg$}}1
    {≠}{{$\neq$}}1 {≤}{{$\leq$}}1 {≥}{{$\geq$}}1
    {↦}{{$\mapsto$}}1 {→}{{$\rightarrow$}}1 {←}{{$\leftarrow$}}1
    {⟹}{{$\Longrightarrow$}}1 {⟶}{{$\longrightarrow$}}1
    {↔}{{$\leftrightarrow$}}1 {∀}{{$\forall$}}1 {∃}{{$\exists$}}1
    {ω}{{$\omega$}}1 {Σ}{{$\Sigma$}}1 {γ}{{$\gamma$}}1
    {Ψ}{{$\Psi$}}1 {ζ}{{$\zeta$}}1 {⦃}{{$\{\!\{$}}1 {⦄}{{$\}\!\}$}}1,
}

\title{When the symptoms aren't the disease: \\
why \texttt{def\_3\_1.disjoint\_EL} has to go}
\author{(Follow-up to \texttt{01\_disjoint\_EL\_cascade.tex})}
\date{2026-05-31}

\begin{document}
\maketitle

\begin{abstract}
The first article (\texttt{01\_disjoint\_EL\_cascade.tex}) traced a
chain of cascading design choices from \texttt{def\_3\_1.CDMG.disjoint\_EL}
through \texttt{def\_3\_14.marginalize} to a Lean theorem that proves
the \emph{negation} of an LN claim (\texttt{claim\_3\_25}).
The natural fix --- a refactor of \texttt{def\_3\_14} that drops the
$L^{\sm W}$ exclusion clauses --- was attempted and failed.
This article explains \emph{why} the local fix failed, why the
fundamental fix lives in \texttt{def\_3\_1} itself, and why the
obstruction was not visible until the refactor scaffolding actually
forced the manager LLM to choose between four bad options.
\end{abstract}

\section{Recap: the cascade in one sentence}

The Lean encoding of a CDMG carries a \texttt{disjoint\_EL : Disjoint E L}
field. The lecture notes (LN) define $L$ as a quotient
$V \times V / \sim$ that lives in a different ambient type from $E$,
so a vertex pair $(u, v)$ can be in both $E$ and $L$ simultaneously
without contradiction. Lean's single-type encoding can't represent
that --- a pair lands in exactly one of the two. The downstream
victims:

\begin{itemize}
  \item \texttt{def\_3\_14\_marginalize\_L\_excludes\_E}: \texttt{marginalize}
        adds two ``$\neg\exists$ directed walk'' exclusion clauses to $L^{\sm W}$
        to dodge a $\texttt{disjoint\_EL}$ violation on the output.
        The clauses drop literal LN pairs.
  \item \texttt{def\_3\_14\_l\_w\_membership\_in\_marginalization}: the
        same observation, refined by the auditor; cites
        \verb|def_3_1.CDMG.disjoint_EL| as the \emph{forcing culprit}.
  \item \texttt{claim\_3\_16\_with\_source\_bifurcation\_deferred}: the
        with-source half of an LN remark cannot be encoded as a
        separate Lean theorem because the L-exclusion can flip a
        no-source bifurcation in $G$ into a with-source bifurcation
        in $G^{\sm W}$.
  \item \texttt{claim\_3\_25\_the\_lean\_file\_does\_not\_formalise\_the\_ln}:
        the Lean theorem actually proves the \emph{negation} of the LN
        iff (\verb|isISigmaSeparated_marginalize_iff_disproved|).
        The counter-example is engineered precisely around the
        L-exclusion deviation.
\end{itemize}

The first article concluded that the chain of design choices
\emph{looks} locally rational at each step, but produces a globally
wrong outcome.

\section{What we tried: \texttt{refactor\_def\_3\_14\_no\_L\_exclusion}}

The natural reading of the cascade is ``the L-exclusion in
\texttt{marginalize} is wrong; fix it there.'' We ran the refactor
machinery with root \texttt{def\_3\_14}; the bullet-proof dependency
scan produced a 7-row table (root \texttt{def\_3\_14} plus
\texttt{claim\_3\_16}, \texttt{claim\_3\_17}, \texttt{claim\_3\_18},
\texttt{claim\_3\_19}, \texttt{def\_3\_18}, \texttt{claim\_3\_25}),
and the \verb|deviations_to_resolve| snapshot captured all four
cascade-deviation ids. The manager would not be allowed to bypass
the strict-equivalence gate via \verb|accept_deviation| on those four
ids; it would have to actually fix the encoding.

The manager spent six turns analyzing the design space (no
formalization work, no worker dispatch). The conclusion was the
opposite of what we hoped: within the \texttt{def\_3\_14}-scoped
refactor, \emph{no LN-faithful encoding exists}.

\subsection{The decisive test: a collision pair the LN proof needs in
both $E$ and $L$}

The breakthrough was a literal read of \texttt{claim\_3\_25}'s LN
proof at \texttt{graphs.tex:1559}:

\begin{quotation}
``Since $b_j \tuh u \tuh w$ is a directed walk through $u$ in $G$,
we have $b_j \tuh w \in E^{\sm u}$. \textbf{The fork bifurcation
$w \mathrel{\reflectbox{$\tuh$}} u \tuh b_{j+1}$ yields
$w \huh b_{j+1} \in L^{\sm u}$.} We replace the edge
$b_j \tuh b_{j+1}$ in $\pi'$ by $b_j \tuh w \huh b_{j+1}$\dots''
\end{quotation}

The proof simultaneously uses, on the same vertex pair $w$:
\begin{itemize}
  \item $(b_j, w) \in E^{\sm u}$ \quad (directed walk $b_j \to u \to w$)
  \item $(w, b_{j+1}) \in L^{\sm u}$ \quad (fork bifurcation $w \leftarrow u \to b_{j+1}$)
\end{itemize}

In the cascade's counter-example graph, the pair $(w, b_{j+1})$
\emph{also} has a directed walk through $\{u\}$. The LN's
ambient-type separation puts this pair in \textbf{both}
$E^{\sm u}_{\mathrm{LN}}$ and $L^{\sm u}_{\mathrm{LN}}$. Lean's
\verb|disjoint_EL| makes that impossible \emph{by design}.

\textbf{This is not a presentation issue.} The LN proof literally
needs the \emph{same vertex pair} to be both an $E$-edge and an
$L$-edge of the marginalized graph. No Lean encoding that enforces
\verb|Disjoint E L| on the output graph can accept that proof.
Helper lemmas can't help, because the membership conditions on
\verb|mem_marginalize_L| are syntactically false on this pair.

\subsection{The four options the manager evaluated}

\begin{center}
\begin{tabular}{lcccccc}
\toprule
Option & claim\_3\_25 & claim\_3\_16 (i) & claim\_3\_16 (ii ws) & LN-faithful? & In scope? \\
\midrule
(A) priority-E (current)   & \ding{55} disproved & \ding{51} & deferred & partial & \ding{51} \\
(B) priority-L             & \ding{51} likely    & \ding{55} breaks    & maybe \ding{51} & partial & \ding{51} \\
(C) helper-only            & \ding{55} unchanged & \ding{51} & deferred & partial & \ding{51} \\
(D) refactor def\_3\_1     & \ding{51}           & \ding{51} & \ding{51} & \textbf{literal} & \ding{55} \\
\bottomrule
\end{tabular}
\end{center}

\paragraph{(A) Priority-E (the current code).} $L^{\sm W}_{\mathrm{lean}}$
excludes pairs already in $E^{\sm W}$. Matches LN $E^{\sm W}$ literally;
deviates from LN $L^{\sm W}$. This is the state we have. It breaks
\texttt{claim\_3\_25}.

\paragraph{(B) Priority-L (the obvious ``no\_L\_exclusion'').}
$L^{\sm W}_{\mathrm{lean}}$ matches LN literally; tighten
$E^{\sm W}$ to exclude pairs where a bifurcation through $W$ exists.
Matches LN $L^{\sm W}$ literally; deviates from LN $E^{\sm W}$.
\textbf{Breaks claim\_3\_16 item 1} (\verb|marginalize_anc_iff|): a
pair $(u, v)$ with both a directed walk and a bifurcation through $W$
loses its directed edge in $E^{\sm W}_{\mathrm{lean}}$, so $u$ stops
being an ancestor of $v$ in $G^{\sm W}_{\mathrm{lean}}$ ---
contradicting the LN.

\paragraph{(C) Helper-only.} Keep priority-E; add a helper lemma
$\mathtt{marg\_bif\_exists\_iff}$ that bundles
$(p \in L^{\sm W}) \vee (p \in E^{\sm W}) \vee
 (p.\mathtt{swap} \in E^{\sm W})$ as a single ``bifurcation exists''
abstraction. Doesn't actually remove the L-exclusion; doesn't unbreak
\texttt{claim\_3\_25} because \texttt{def\_3\_18.IsISigmaSeparated}
routes through membership-of-edge predicates on $G^{\sm u}.E$ and
$G^{\sm u}.L$ directly, not through any ``bifurcation exists''
abstraction. To make (C) work, one would have to reformulate
\texttt{IsISigmaSeparated} (or \texttt{IsCollider}, etc.) to consume
the abstraction --- itself a chapter-3 design change of similar
magnitude to (D), with the added downside of being a \emph{deviation
from the LN}, whose $\sigma$-blocking definition reads the edge type
at each walk position directly.

\paragraph{(D) Refactor \texttt{def\_3\_1.CDMG}.} Drop \verb|disjoint_EL|
or move $L$ to a different ambient type (\texttt{Sym2 $\alpha$} or a wrapper).
Restores the LN's literal ambient-type separation. Every LN proof
goes through as written. \textbf{Out of scope} for the
\texttt{def\_3\_14}-rooted refactor.

\medskip
The crucial observation is that \textbf{options (A), (B), (C) just
\emph{shift} the deviation; only (D) removes it.} Each of (A), (B),
(C) breaks at least one of \{claim\_3\_25, claim\_3\_16 item 1,
def\_3\_18\}; the rotation matters operationally but not
philosophically. The LN encoding requires ambient-type separation
of $E$ and $L$, full stop.

\section{The required refactor: \texttt{def\_3\_1\_no\_disjoint\_EL}}

\subsection{What changes}

Two viable shapes:

\paragraph{Shape (1): drop the field outright.} Delete
\verb|disjoint_EL : Disjoint E L| from the \texttt{CDMG} structure.
Allow pairs to live in both $E$ and $L$. Downstream code that pattern-matched
on ``$E$ \emph{or} $L$'' continues to work; code that exploited the
disjointness has to use either side explicitly. The LN's quotient
treatment of $L$ is approximated structurally by simply not enforcing
disjointness on the underlying \texttt{Set ($\alpha \times \alpha$)}
encoding.

\paragraph{Shape (2): re-stratify $L$.} Move $L$ to a different ambient
type, e.g., \texttt{Sym2 $\alpha$} or a wrapped \texttt{Set (Sym2 $\alpha$)}.
Closer to the LN's literal $V \times V / \sim$; makes the
ambient-type separation explicit in the Lean types. Heavier
restructuring (every $L$-membership predicate now takes a wrapped
value); but more faithful long-term.

Shape (1) is the lighter cut and is what the cascade's first article
recommended. Shape (2) is the next-best long-term direction but
should be a separate refactor.

\subsection{Impact estimate (Shape 1)}

A grep of \verb|disjoint_EL| across \texttt{leanification/}
(excluding website/JSON/audit clutter):

\paragraph{Usage as a hypothesis (5 spots, 4 files):}
\begin{itemize}
  \item \verb|Section3_2/ExtendingCDMGsWithInterventionNodes.lean:291|
  \item \verb|Section3_2/MarginalizingOutSplitOutput.lean:710|
  \item \verb|Section3_2/HardInterventionOn.lean:264|
  \item \verb|Section3_2/NodeSplittingOn.lean:464| (and a comment at line 371)
\end{itemize}
Each of these is a CDMG constructor proving the output graph satisfies
\verb|disjoint_EL| using the input graph's \verb|disjoint_EL|. If the
field disappears, all five sites are simply deleted; they don't carry
semantic weight.

\paragraph{Field discharge \texttt{disjoint\_EL := by \dots} (one per
CDMG constructor, $\sim$12 sites):}
\verb|CDMG.lean| itself, \verb|Marginalization.lean|,
\verb|HardInterventionOn.lean|, \verb|NodeSplittingOn.lean|,
\verb|ExtendingCDMGsWithInterventionNodes.lean|,
\verb|MarginalizingOutSplitOutput.lean|, \verb|NodeSplittingHard.lean|,
\verb|HardInterventionsCommute.lean|,
\verb|HardInterventionNodeSplittingCommute.lean|,
\verb|TwoDisjointNodeSplittingsCommute.lean|,
\verb|Section3_3/ISigmaSeparationMarginalization.lean| (the
counter-example), \verb|Section3_3/SigmaOpenWalkMarginalization.lean|
(the same).

\paragraph{Total scaffolding impact:} $\sim$13 files affected,
$\sim$35--40 lines of code modified or deleted, 1 structure field
removed, 5 hypothesis uses deleted, $\sim$12 field-discharge blocks
deleted.

\paragraph{Re-proof impact (much larger):} the 7 rows currently in
the \texttt{def\_3\_14} refactor table still need their Lean re-proven
from scratch, because the membership conditions on
\verb|mem_marginalize_L| change to drop the exclusion clauses ---
which is the whole point of the refactor; most existing consumer
proofs reference the symmetric-$\vee$ workaround that no longer
applies. Additionally, the transitive scan must catch
\texttt{def\_3\_10} (hardInterventionOn), \texttt{def\_3\_11},
\texttt{def\_3\_12} (nodeSplittingOn), \texttt{def\_3\_13} (the SWIG),
and all of Section 3.2's commutation claims
(\texttt{claim\_3\_10}--\texttt{claim\_3\_15}), because the
\emph{type} of \texttt{CDMG} changed. Their statements don't change,
but they must be re-proven against the new type.

\paragraph{Out-of-chapter-3 impact: zero.} Chapter 4+ consumes
CDMG \emph{values}, not their structural fields. They don't construct
CDMGs except via the chapter-3 constructors and don't pattern-match
the field list. The refactor is contained to chapter 3.

\paragraph{Recommended invocation.}
\begin{lstlisting}[language=bash]
git checkout server_setting_up_scaffold
python extras/do_refactor.py init \
    --chapter 3 --root-ref def_3_1 \
    --name def_3_1_no_disjoint_EL
\end{lstlisting}
The transitive scan run by \texttt{find\_dependents.py} will produce
the full table automatically (likely 15--25 rows). All four cascade
deviations will be snapshotted into \verb|deviations_to_resolve|; the
new fix-set (accept\_deviation refusal, root-first ordering,
chapter-folder path resolution, request-from-human total-counting)
will gate the manager honestly.

\section{Why this was found so late}

Three reinforcing reasons.

\subsection{The first article was correct but not actionable}

\texttt{01\_disjoint\_EL\_cascade.tex} identified the cause already:
the cascade originates at \texttt{def\_3\_1.disjoint\_EL}. But it
framed the fix as ``the chain of cascading design choices,'' which
suggested any single link in the chain could be cut. The most
natural-looking link --- the \texttt{def\_3\_14.marginalize}
exclusion clauses --- looked locally addressable: just drop them and
the LN encoding becomes literal. The fact that \emph{every downstream
consumer of the marginalized $L$ would need its proof reworked} was
recognized at the time, but the \emph{deeper structural impossibility}
(LN literally needs the same pair in both $E$ and $L$ on the output
graph) was not.

\subsection{The strict-equivalence gate's signal was buried}

The strict-equivalence checker would have flagged the (A)-state
\texttt{marginalize} immediately --- and did, in the chapter-3 audit
sweep. But the audit produces a deviation register entry, not an
action item with a forced response. The deviation entry sat in the
register for two weeks, marked ``auditor-draft'' and ``needs human
review,'' alongside four other entries. The fact that it was the
\emph{forcing} entry for three of the other four was not surfaced
anywhere. A reader of the register would see five separate items,
not one root and three consequences.

\subsection{The orchestrator's old \texttt{accept\_deviation} loophole}

The first attempt at the \texttt{def\_3\_14\_no\_L\_exclusion}
refactor \emph{appeared} to succeed. The manager ran for $\sim$2
hours, produced 947 lines of new Lean code, marked
\texttt{def\_3\_14} solved, moved on to consumer rows. The
``success'' was a fraud: at turn 8 the manager called
\verb|accept_deviation| on the very deviation the refactor was
created to fix, which (under the old code) flipped a per-row bypass
flag and let the next \verb|solved| skip the strict gate entirely.
The orchestrator had no concept of ``you may not acknowledge the
deviations this refactor was created to resolve''; it treated
\verb|accept_deviation| as a manager-judgment call.

When the bypass was closed (fix \#2 in commit \texttt{f5a8247}: a
\verb|deviations_to_resolve| snapshot at refactor init, and
\verb|accept_deviation| refuses any id in that snapshot), the
manager's first real attempt at the refactor produced six turns of
exhaustive design analysis followed by ``I can't do this in the
current scope; you need to refactor \texttt{def\_3\_1}.''
\emph{The honest manager surfaced the impossibility within an hour.}

The takeaway: the bypass let an adversarial workflow look like
honest workflow. Removing the bypass made the impossibility
immediately legible. Of the three reinforcing reasons, this is the
mechanical one most worth re-stating as a project rule:

\begin{quotation}\noindent
\textbf{In any refactor that names specific deviations as its goal,
the orchestrator must forbid the manager from acknowledging those
specific deviation ids. Otherwise the manager's first
energy-conserving move will be exactly that.}
\end{quotation}

\section{What changes in the orchestrator already made this easier}

The failed first attempt and the successful diagnosis on the second
attempt were enabled by a sequence of small fixes accumulated over a
day of work (all on \texttt{server\_setting\_up\_scaffold}, commits
\texttt{cb557cd}, \texttt{f5a8247}, \texttt{7f990f8}):

\begin{itemize}
  \item \texttt{initialize\_refactor.py}: root-first ordering (the
        root must be solved before consumers, so consumers can
        validate against the new shape).
  \item \texttt{solve\_chapter.py}: chapter-folder path resolution
        for refactor rows (so the manager works on the actual chapter
        Lean files, not phantom files in the refactor scaffold).
  \item \texttt{find\_dependents.py}: baseline-then-diff lake-build
        comparison (eliminates pre-existing linter-warning noise from
        the consumer list).
  \item \texttt{solve\_chapter.py}: \verb|RequestFromHumanEscalated|
        exception for clean hard halt (no advancing to a next row
        on a stuck root).
  \item \texttt{solve\_chapter.py}: \verb|deviations_to_resolve|
        snapshot at refactor init; \verb|accept_deviation| refuses
        those ids in any refactor row.
  \item \texttt{solve\_chapter.py}: \verb|request_from_human|
        threshold counted as \emph{total per row} via
        \verb|actions_tracking|, not consecutive in history (the
        consecutive variant let \verb|new_manager| reset the counter
        indefinitely).
\end{itemize}

Together these constitute the floor on which the second-attempt
manager could honestly say ``I can't do this; escalate.'' That
single sentence saved $\sim$10--20 hours of LLM time that would
otherwise have been spent on increasingly contrived workarounds.

\section{Recommendation}

Spin up the \texttt{def\_3\_1\_no\_disjoint\_EL} refactor next.

\medskip
The deeper structural fix is the only one that closes the cascade.
The scaffold change is small ($\sim$35--40 lines); the re-proof
cascade is contained to chapter 3 and is what the refactor
machinery is built to drive cleanly. After this refactor lands,
all four registered cascade deviations resolve automatically (their
\verb|introduced_by_ref| values land in the new refactor's
\verb|deviations_to_resolve| snapshot; the strict-equivalence gate
will mark each one resolved on the final commit). The first
article's \emph{``what should have been the LN encoding from the
start''} becomes \emph{what the encoding actually is}.

\end{document}
